% ======================================================================
%  METHODOLOGY
% ======================================================================
\section{Methodology}\label{sec:method}

\subsection{Cross-Exchange Spread and Z-Score Target}

Let $P_i^c(t)$ denote the mid-price of asset $c$ on exchange $i$ at snapshot $t$. For an exchange pair $(i, j)$, the cross-exchange spread in basis points is:
\begin{equation}
    S_{ij}^c(t) = \frac{P_i^c(t) - P_j^c(t)}{P_i^c(t)} \times 10{,}000
\end{equation}

Unlike traditional pairs trading where the spread between different assets may diverge permanently as fundamentals change, the cross-exchange spread for the same asset is anchored by the Law of One Price: absent frictions, $S_{ij}^c(t) \to 0$. In practice, fees, withdrawal delays, and matching-engine latency allow temporary price deviations for the same coin across different exchanges.

We normalize the spread into a rolling z-score:
\begin{equation}
    z_t = \frac{S_t - \hat{\mu}_w}{\hat{\sigma}_w}
\end{equation}
where $\hat{\mu}_w$ and $\hat{\sigma}_w$ are the estimated rolling mean and standard deviation over a window of $w$ snapshots, with a minimum period requirement before the z-score is defined. Model outputs $\hat{z}$ are therefore in units of this trailing $\hat{\sigma}_w$: $|\hat{z}|=1$ is a one-$\sigma$ forecasted dislocation. We use $w=300$ ($\text{MIN\_PERIODS}=90$); larger $w$ adds context but lengthens warmup (Section~\ref{sec:ablation}).

\subsection{Supervised Prediction Target}

The model's target is the \emph{future} z-score at horizon $H$:
\begin{equation}
    y_t = z_{t+H}
\end{equation}
computed via $\texttt{groupby(coin, pair).zscore.shift(-H)}$ on the training panel. This is a regression target, not a binary classification. The model predicts both the direction and magnitude of the z-score $H$ snapshots ahead, enabling a confidence filter on $|\hat{z}|$.

For the prediction model used in the paper-trading campaigns, $H = 1$ (predict the next snapshot's z-score), $w = 300$ snapshots, and $\text{MIN\_PERIODS} = 90$. Window length and confidence cutoff ($\tau$) selection is described in Section~\ref{sec:ablation}.

\subsection{Feature Engineering}\label{sec:features}

The model uses a 68-dimensional feature matrix per (coin, exchange-pair, snapshot), built from four synchronized live signal families: (i)~lagged cross-exchange spreads and rolling z-scores, with short momentum and acceleration of those series; (ii)~ticker top-of-book mids, bid--ask widths, and volumes; (iii)~L2 orderbook imbalance over the top~20 levels on each venue,
\begin{equation}
    \texttt{imbalance}_i = \frac{V_{\text{bid},i} - V_{\text{ask},i}}{V_{\text{bid},i} + V_{\text{ask},i}},
\end{equation}
with short lags of that ratio; and (iv)~recent trade-flow intensity and buy/sell imbalance, plus cross-venue aggregates (mid dispersion, BA/OB variability, and price ranks). $\texttt{coin}$ and $\texttt{pair}$ enter as LightGBM native categoricals so one panel model can specialize by asset and venue pair. Exact column names and construction details ship with the public dataset release (\datasetref); Section~\ref{sec:ablation} ablates the engineered spread/$z$ trajectory against raw live columns.

\subsection{Model: LightGBM Gradient Boosting}\label{sec:lgbm}

We train a single LightGBM~\cite{ke2017lightgbm} regression model on the pooled panel of all (coin, exchange-pair, snapshot) observations. The objective is squared error minimization on $y_t = z_{t+H}$.

We deliberately choose tree ensembles over the DNN/LSTM architectures common in prior spread-forecasting work: the feature matrix is a heterogeneous tabular panel (lags, microstructure, and native categoricals) for which tree models are a strong default inductive bias. As a classical Random Forest baseline on the same train/test feature matrix, we also fit a Random Forest regressor~\cite{breiman2001random} (regularized leaf size analogous to LightGBM's minimum child samples). LightGBM is the model used for validation-campaign evaluation; the Random Forest control is scored on the test set (Section~\ref{sec:results}).

\subsection{Trading Policy}\label{sec:policy}

At each live snapshot, the model generates a prediction $\hat{z}_{t+1}$ for every (coin, exchange-pair) observation. The trading policy is:

\begin{enumerate}
    \item Entry: open a position if $|\hat{z}_{t+1}| \geq \tau$ and the (coin, pair) slot is not already occupied. The paper protocol uses $\tau{=}0.9$.
    \item Direction: $\text{sign}(\hat{z}_{t+1})$. The model predicts the sign of the spread z-score at $t+1$ and bets accordingly.
    \item Capacity: at most $\text{max\_open} = 50$ positions open simultaneously, enforced by a first-come queue.
    \item Settlement: each position settles at $t+H$ (with $H=1$). The proxy profit is $\text{pnl\_proxy} = \text{direction} \times z_{t+H}$.
\end{enumerate}

The settlement proxy measures whether the z-score moved in the predicted direction and by how much, in z-score units. It is not a dollar profit metric and does not account for transaction fees or slippage.


\subsection{Evaluation Metrics}

We report both forecast quality and trading performance metrics:

\begin{itemize}
    \item \textbf{R$^2$}: fraction of variance in the realized $z_{t+H}$ explained by the model's predictions.
    \item \textbf{Directional accuracy (DirAcc)}: fraction of predictions where $\text{sign}(\hat{z}) = \text{sign}(z_{t+H})$.
    \item \textbf{Mean pnl\_proxy}: average $\text{direction} \times z_{t+H}$ across settled trades.
    \item \textbf{Per-trade Sharpe}: $\text{mean}(\text{pnl\_proxy}) / \text{std}(\text{pnl\_proxy})$ on the filtered trade set (Rf~$=0$), used to select $\tau$ in Section~\ref{sec:ablation}.
    \item \textbf{Hourly portfolio Sharpe ratio}: $\text{mean}(\Delta\text{equity}_h) / \text{std}(\Delta\text{equity}_h)$ where hourly equity includes realized closed-trade profits plus mark-to-market valuation of open positions as $\text{direction} \times z_t$. Risk-free rate is zero because the z-unit proxy has no currency denomination and represents forecast quality rather than economic returns.
\end{itemize}

All metrics are computed with a matched \textbf{naive persistence baseline} ($\hat{z}_{t+H} \leftarrow z_t$) scored on the identical row set.
